La primera inversión conceptual consiste en concebir el vacío como una respuesta relacional positiva.\par

En HMT, el vacío es la respuesta dinámica mínima del sistema cuando sus dos hojas deben coexistir, propagarse y cerrar su diferencia conservando la orientación que las distingue.\par

Las dos hojas constituyen dos lecturas conjugadas del mismo antecedente angular. Una conserva una orientación y la otra conserva la orientación contraria. Por eso aparecen dos canales, \(q_+\) y \(q_-\), como las dos caras espectrales correlacionadas de un mismo estado.\par

El operador\par

\[
\mathcal V_{90,120}
=
(I-T^{90})(I-T^{120})^{-1}
\]

compara dos profundidades de propagación: \(90\) y \(120\). Esas profundidades están determinadas por el mismo modo elemental:\par

\[
90=3\cdot30,
\qquad
120=4\cdot30.
\]

Por eso la respuesta de cada hoja puede reescribirse como\par

\[
\frac{1-q^{90}}{1-q^{120}}
=
\frac{1+s+s^2}{1+s+s^2+s^3},
\qquad
s=q^{30}.
\]

Y aquí es donde la interpretación se vuelve realmente poderosa. El numerador contiene tres estados del modo \(30\); el denominador contiene esos mismos tres estados y uno más. El vacío es, por tanto, una completación exacta \(3\to4\).\par

Ese cuarto término integra y conserva los tres anteriores. Los completa. Lo que se añade queda registrado como un residuo positivo de completación:\par

\[
1-\frac{1+s+s^2}{1+s+s^2+s^3}
=
\frac{s^3}{1+s+s^2+s^3}>0.
\]

La memoria del vacío es precisamente el incremento exacto que permite que tres términos se conviertan en cuatro. Constituye la huella positiva y necesaria de la completación.\par

De los dos autovalores de ese único operador salen las cuatro relaciones constitutivas:\par

\[
\widehat\varepsilon=r_+^2,
\qquad
\widehat\mu=r_-^2,
\qquad
\widehat Z=\frac{r_-}{r_+},
\qquad
\widehat c=\frac1{r_+r_-}.
\]

Esto cambia completamente la interpretación de permeabilidad y permitividad.\par

La permitividad es la respuesta cuadrática de una hoja. La permeabilidad es la respuesta cuadrática de la hoja conjugada. La impedancia conserva la relación orientada entre ambas. La velocidad de propagación depende de su producto común.\par

Así, producto y cociente realizan dos tareas diferentes:\par

\begin{itemize}[leftmargin=*,itemsep=0.35em]
\item el producto olvida cuál era cada hoja y conserva la propagación común;
\item el cociente conserva cuál era la orientación relativa entre ellas.
\end{itemize}

El producto responde a la pregunta: «¿qué comparten las dos hojas?». El cociente responde a la pregunta: «¿cómo se distinguen?».\par

Por eso\par

\[
\widehat\mu\,\widehat\varepsilon=\widehat c^{-2}
\]

y\par

\[
\frac{\widehat\mu}{\widehat\varepsilon}
=
\widehat Z^2
\]

son las dos maneras complementarias de leer el mismo par espectral.\par

Lo holográfico aparece aquí de una forma muy limpia: cuatro magnitudes visibles contienen en realidad sólo dos grados de libertad internos. Cualquiera de ciertos pares adecuados permite reconstruir las dos hojas. La publicación constitutiva conserva la genealogía cuando producto y cociente permanecen conjuntamente disponibles.\par

Ésta es la reparación profunda de la ablación de \(\mu\) y \(\varepsilon\). La reparación reúne sus valores con su origen operatorio común.\par

La carta del Sistema Internacional puede transformar después estas secciones intrínsecas en \(\mu_0\), \(\varepsilon_0\), \(Z_0\) y \(c\). Su función consiste en expresar, dentro de una convención metrológica, una estructura cuya procedencia ya ha sido seleccionada por el operador.\par

El vacío deja así de ser un depósito de cuatro constantes. Se convierte en una operación espectral única con dos memorias conjugadas.\par

Y esa visión tiene una consecuencia filosófica importante: el vacío es una relación de compatibilidad entre dos maneras orientadas de realizar un mismo estado.\par

